\clearpage
\SourceLeaf{362}%
\pdfbookmark[0]{Anzeigen. Mollweide: Archimedes. 1808.}{anzeigen-mollweide-1808}
\begin{center}
\rule{.85\textwidth}{.3pt}\par
Göttingische gelehrte Anzeigen.\quad 9. Januar 1808.\par
\rule{.85\textwidth}{.3pt}
\end{center}
\hspace*{\parindent}\textsc{Archimedes} bestimmte bekanntlich in seiner Schrift \textit{Circuli dimensio} die Grenzen für den Umfang des Kreises, indem er ihn zwischen den Umfang eines umbeschriebenen und den eines einbeschriebenen 96-Ecks einschloss. Diese Zahlen berechnet er, oder genauer gesagt bestimmt er eine größere Zahl als die erste und eine kleinere als die zweite, indem er schrittweise vom Sechseck zum Zwölfeck, von diesem zum 24-Eck und so weiter fortschreitet. Für beide 96-Ecke geht er daher, in unserer Ausdrucksweise, von einem Näherungswert der irrationalen Größe \(\BareRoot{3}\) aus. Der eine, nämlich \(\frac{265}{153}\), ist etwas zu klein, der andere, \(\frac{1351}{780}\). etwas zu groß. Jener wird bei den umbeschriebenen, dieser bei den einbeschriebenen Vielecken verwendet. Bei genauerem Hinsehen zeigt sich, dass diese Näherungswerte in der Folge \(\frac21,\frac53,\frac74,\frac{19}{11},\frac{26}{15}\) und so weiter vorkommen. Ihre Glieder sind abwechselnd größer und kleiner als \(\BareRoot{3}\), und jedes weicht weniger davon ab als irgendein anderer durch kleinere Zahlen ausgedrückter Bruch. Der Bruch \(\frac{265}{153}\) ist nämlich das achte und \(\frac{1351}{780}\) das elfte Glied der Folge. Es scheint also, dass \textsc{Archimedes} diese Näherungswerte nach einer Methode gefunden hat, nicht durch Zufall. Da er sich aber selbst überhaupt nicht dazu äußert, wie er zu ihnen gelangt ist, und sich auch sonst kein Hinweis darauf findet, dass die Alten unsere Methoden zur Lösung solcher Aufgaben kannten, ist hier Raum für Vermutungen. Professor \textsc{Mollweide} in Halle hat die Frage in einer kurzen Abhandlung untersucht, die er kürzlich der Königlichen Sozietät, deren Korrespondent er ist, unter folgendem Titel eingesandt hat:
\begin{center}\textit{De methodo ab Archimede adhibita ad rationem, in qua inter se sunt latus trianguli aequilateri et radius circuli circumscripti, numeris veritati proxime exprimendam}\end{center}
\noindent Er gibt darin ein Verfahren an, das dem Stand der Arithmetik der Alten entspricht und daher vielleicht genau dasjenige ist, das \textsc{Archimedes} verwendet hat. Herr M. bezeichnet die Seite des Dreiecks mit \(AC\) und den Radius des %
